\documentclass[11pt,a4paper]{article}
\usepackage[utf8]{inputenc}
\usepackage[T1]{fontenc}
\usepackage[norsk]{babel}
\usepackage{tgtermes}
\usepackage{tgheros}
\renewcommand{\ttdefault}{txtt}
\usepackage{amsmath,amssymb}
\usepackage[a4paper,left=2.45cm,right=2.6cm,top=2.3cm,bottom=2.2cm,headsep=0.9cm]{geometry}
\usepackage{fancyhdr}
\usepackage{lastpage}
\usepackage{xcolor}
\usepackage{tikz}
\usepackage{pgfplots}
\pgfplotsset{compat=1.18}
\usepackage{booktabs}
\usepackage{array}
\usepackage{enumitem}
\usepackage{listings}
\usepackage{needspace}
\usepackage{multirow}
\usetikzlibrary{positioning,arrows.meta,calc}
\definecolor{insgrey}{gray}{0.42}
\definecolor{radiogrey}{gray}{0.70}
\setlength{\parindent}{0pt}
\setlength{\parskip}{0.45em}
\setlist[itemize]{itemsep=0pt,topsep=2pt,label=\textbullet,font=\sffamily}
\lstset{basicstyle=\ttfamily\small,columns=fullflexible,keepspaces=true,
  showstringspaces=false,upquote=true,aboveskip=0.6em,belowskip=0.6em,
  literate={æ}{{\ae}}1 {ø}{{\o}}1 {å}{{\aa}}1 {Æ}{{\AE}}1 {Ø}{{\O}}1 {Å}{{\AA}}1}
\pagestyle{fancy}
\fancyhf{}
\renewcommand{\headrulewidth}{0pt}
\fancyhead[R]{\fontsize{7.5}{9}\selectfont\sffamily\color{insgrey}TDT4172 - Introduksjon til maskinlæring - Øvingssett 8}
\fancyfoot[R]{\fontsize{7}{9}\selectfont\sffamily\color{insgrey}\thepage/\pageref{LastPage}}
\newcommand{\delstart}[1]{\clearpage
  \noindent\hspace*{-1.55em}\makebox[1.55em][l]{\sffamily\bfseries\fontsize{12}{14}\selectfont\raisebox{0.35em}{#1}}%
  {\sffamily\bfseries\fontsize{18}{22}\selectfont Part #1}\par\vspace{1.0em}}
\newcommand{\sporsmal}[1]{\par\Needspace{9\baselineskip}\vspace{1.9em}%
  \noindent{\bfseries Spørsmål #1}\par\vspace{-0.45em}%
  \noindent{\bfseries -{}-{}-{}-{}-{}-{}-{}-{}-{}-{}-{}-{}-{}-{}-{}-{}-}\par\vspace{-0.15em}}
\newcommand{\velg}[1]{\par\Needspace{6\baselineskip}\noindent{\bfseries #1}\par\vspace{0.3em}}
\newcommand{\radio}[1]{\tikz[baseline=-0.62ex]{\draw[radiogrey,line width=0.45pt,fill=gray!5] (0,0) circle (0.56em);
  \node[inner sep=0pt,font=\sffamily\fontsize{5.5}{6}\selectfont,text=insgrey] at (0,0) {#1};}}
\newenvironment{alternativer}{\begin{list}{}{\setlength{\leftmargin}{2.35em}\setlength{\labelwidth}{1.4em}%
  \setlength{\labelsep}{0.6em}\setlength{\itemsep}{0.85em}\setlength{\topsep}{0.45em}\setlength{\parsep}{0pt}}%
  \sffamily\fontsize{10}{12.6}\selectfont}{\end{list}}
\newcommand{\alt}[2]{\item[\radio{#1}]#2}
\newcommand{\makspoeng}{\par\vspace{2.2em}\hfill{\sffamily\fontsize{8.5}{10}\selectfont Maks poeng: 10}\hspace*{0.5em}\par}
\begin{document}

\noindent\hspace*{-1.55em}\makebox[1.55em][l]{\raisebox{0.3em}{\sffamily\bfseries\fontsize{11}{12}\selectfont i}}%
{\sffamily\bfseries\fontsize{18}{22}\selectfont Frontpage}\par\vspace{1.1em}

{\bfseries Eksamensoppgave i TDT4172 - Introduksjon til maskinlæring}\\
{\bfseries Øvingssett:} 8 av 10 (uoffisielt øvingssett)\\
{\bfseries Tid:} 4 timer

\vspace{0.6em}
{\bfseries Faglig kontakt:} --- (øvingssett, ikke en offisiell eksamen)\\
{\bfseries Møter i eksamenslokalet:} Nei

\vspace{0.6em}
{\bfseries Hjelpemiddelkode/Tillatte hjelpemidler:} D / Ingen trykte eller håndskrevne hjelpemidler tillatt.
Bestemt, enkel kalkulator tillatt.

\vspace{0.6em}
{\bfseries\underline{ANNEN INFORMASJON:}}

\vspace{0.4em}
{\bfseries Les oppgavene nøye og gjør dine egne antagelser.} Presiser i besvarelsen hvilke forutsetninger du
har lagt til grunn i tolkning/avgrensing av oppgaven.

Dette øvingssettet er laget for egentrening, med samme format og vanskelighetsgrad som de
siste ordinære eksamenene i emnet. Løsningsforslaget står på siste side.

\vspace{1.6em}
{\bfseries\underline{FAGSPESIFIKK INFORMASJON}}

\vspace{0.4em}
{\bfseries Ingen papirhåndtegninger:}\\
Denne eksamenen tillater ikke bruk av håndtegninger.

\vspace{1.4em}
{\bfseries Vekting av oppgavene}:\\
Denne eksamenen inneholder flervalgsoppgaver organisert i ti grupper på fire spørsmål hver. Hver
gruppe vurderes separat etter følgende regler:

\vspace{0.6em}
For hver av de fire spørsmålene i gruppen får du følgende poeng
\begin{itemize}[leftmargin=2.2em,label=\textperiodcentered,itemsep=0pt]
\item {\bfseries Riktig svar:} +2,5 poeng
\item {\bfseries Feil svar (et alternativ er valgt, men ikke korrekt):} $-$0,83 poeng
\item {\bfseries Ikke besvart:} 0 poeng
\end{itemize}
{\bfseries Merk} at dersom du først har valgt et svaralternativ, kan du {\bfseries ikke velge det bort} igjen. Da må du
svare på det spørsmålet. Med andre ord mister du muligheten til å la et spørsmål stå blankt så fort
du har trykket på et alternativ.

\vspace{0.6em}
Poengsummen for hver gruppe er {\bfseries summen} av poengene for de fire enkeltspørsmålene.\\
For {\bfseries alle spørsmålene} gjelder en {\bfseries null-grense}: dersom summen for denne gruppen blir negativ,
settes verdien til {\bfseries 0}. Den totale poengsummen på eksamen er summen av poengene fra alle
gruppene.

\delstart{1}
\sporsmal{1}
Et nevralt nettverk klassifiserer bilder i tre klasser og har \textbf{softmax} i output-laget. Det trenes med
kategorisk cross-entropy $L(\hat{\mathbf{y}}, \mathbf{y}) = -\sum_i y_i \ln \hat{y}_i$, der $\mathbf{y}$ er
one-hot-enkodet og $k$ er den sanne klassen.

Hvilket utsagn er riktig?
\velg{Velg ett alternativ:}
\begin{alternativer}
\alt{A}{Tapet blir 0 så snart den sanne klassen har høyest $\hat{y}_i$, selv om $\hat{y}_k$ er mindre enn 1}
\alt{B}{Softmax brukes fordi den lar flere klasser ha $\hat{y}_i$ nær 1 samtidig, slik at usikre bilder kan få flere labels}
\alt{C}{Tapet summerer $-\ln \hat{y}_i$ over alle klassene, slik at høye $\hat{y}_i$ for gale klasser straffes direkte}
\alt{D}{Tapet blir $-\ln \hat{y}_k$, så det avhenger direkte bare av predikert sannsynlighet for den sanne klassen}
\end{alternativer}
\sporsmal{2}
En logistisk regresjonsmodell predikerer sannsynligheten $p$ for at en kunde kjøper et produkt. Vekten for
featuren \emph{Antall\_besøk} er $w = 0.7$. Husk at $\mathbf{w}^T\mathbf{x} + b = \ln\frac{p}{1-p}$.

Hva skjer når \emph{Antall\_besøk} øker med 1, og alle andre features holdes fast?
\velg{Velg ett alternativ:}
\begin{alternativer}
\alt{A}{Oddsen $p/(1-p)$ multipliseres med $e^{0.7} \approx 2.0$}
\alt{B}{Sannsynligheten $p$ multipliseres med $e^{0.7} \approx 2.0$}
\alt{C}{Sannsynligheten $p$ øker med 0.7}
\alt{D}{Oddsen $p/(1-p)$ øker med 0.7}
\end{alternativer}
\sporsmal{3}
Klassen \texttt{LogisticRegression} fra notatene starter \texttt{fit} med \texttt{self.weights = np.zeros(...)}
og \texttt{self.bias = 0}. Treningssettet har 300 punkter, hvorav 120 har $y = 1$.

Hva blir gjennomsnittlig binary cross-entropy loss (med $\ln$) for prediksjonene som modellen gir
med disse startparametrene, altså før første oppdatering?
\velg{Velg ett alternativ:}
\begin{alternativer}
\alt{A}{$0.25$}
\alt{B}{$0.50$}
\alt{C}{$0.69$}
\alt{D}{$0.67$}
\end{alternativer}
\sporsmal{4}
En logistisk regresjonsmodell med blant annet featuren \emph{Alder} (i år) er trent til konvergens, og
vekten for \emph{Alder} ble $0.06$. En kollega gjør om \emph{Alder} til måneder (ganger alle verdiene med 12),
lar alt annet være likt, og trener på nytt til konvergens. Tapsfunksjonen er konveks.

Hva er det mest sannsynlige resultatet?
\velg{Velg ett alternativ:}
\begin{alternativer}
\alt{A}{Vekten for \emph{Alder} blir omtrent $0.005$, men biasen må endres, så prediksjonene $\hat y$ endres}
\alt{B}{Vekten for \emph{Alder} blir omtrent $0.06$ som før, men prediksjonene $\hat y$ blir høyere}
\alt{C}{Vekten for \emph{Alder} blir omtrent $0.72$, siden verdiene nå er tolv ganger så store}
\alt{D}{Vekten for \emph{Alder} blir omtrent $0.005$, og alle prediksjonene $\hat y$ blir som før}
\end{alternativer}
\makspoeng

\delstart{2}
\sporsmal{1}
To modeller skal oppdage en sjelden feil som finnes i 0.2\,\% av dataene. ROC-kurvene deres er nesten like
(ROC AUC 0.97 og 0.98), mens precision–recall-kurvene viser en tydelig forskjell mellom modellene.

Hva er den beste forklaringen?
\velg{Velg ett alternativ:}
\begin{alternativer}
\alt{A}{Recall avhenger av hvor mange negative det er, mens TPR ikke gjør det, så de to kurvene skiller seg}
\alt{B}{Precision–recall-kurven tar med de svært mange sanne negative (TN), og blir derfor mer følsom}
\alt{C}{FPR deler på det store antallet negative, så mange falske positive gir lav FPR, men lav precision}
\alt{D}{ROC AUC kan typisk ikke brukes ved ubalanserte klasser, siden verdien da alltid havner nær 1}
\end{alternativer}
\sporsmal{2}
To modeller for en alvorlig sykdom er testet på de samme 1000 pasientene, hvorav 100 er syke
(positiv klasse). Modellen skal brukes til screening: alle som flagges, får en enkel og billig oppfølgingstest,
mens syke som ikke flagges, ikke får behandling i tide.
\begin{center}
\begin{tabular}{lcccc}
\toprule
Modell & TP & FN & FP & TN\\
\midrule
A & 85 & 15 & 85 & 815\\
B & 40 & 60 & 40 & 860\\
\bottomrule
\end{tabular}
\end{center}
Hvilken modell bør velges, og hvorfor?
\velg{Velg ett alternativ:}
\begin{alternativer}
\alt{A}{B, siden den har høyere spesifisitet, og det er spesifisiteten som avgjør hvor mange syke som oversees}
\alt{B}{Begge er like gode, siden de gjør like mange feil og har samme precision på disse pasientene}
\alt{C}{B, siden den har høyere precision, så færre friske sendes unødvendig til oppfølgingstest}
\alt{D}{A, siden den fanger opp 85 av de 100 syke, og en unødvendig oppfølgingstest koster lite}
\end{alternativer}
\sporsmal{3}
Modell A har ROC AUC $= 0.82$ og modell B har ROC AUC $= 0.30$ på det samme valideringssettet.

Hvilket utsagn er riktig?
\velg{Velg ett alternativ:}
\begin{alternativer}
\alt{A}{Snur man prediksjonene til B ($p \to 1-p$), får den AUC 0.70, så B inneholder nyttig informasjon}
\alt{B}{B har AUC under 0.5 fordi valideringssettet er ubalansert, noe som alltid trekker AUC under 0.5}
\alt{C}{A rangerer et tilfeldig positivt punkt over et tilfeldig negativt i 82\,\% av tilfellene, og har derfor 82\,\% accuracy}
\alt{D}{B er dårligere enn tilfeldig gjetning, og prediksjonene inneholder derfor ingen nyttig informasjon}
\end{alternativer}
\sporsmal{4}
Et team har 10\,000 transaksjoner, hvorav 1\,\% er svindel. De kopierer først svindelpunktene i
\emph{hele} datasettet til klassene er like store (oversampling), og deler deretter tilfeldig i trening og
test. Modellen får recall 0.99 på testsettet.

Hva er hovedproblemet med fremgangsmåten?
\velg{Velg ett alternativ:}
\begin{alternativer}
\alt{A}{Testsettet får en annen klassefordeling enn treningsdelen, siden kopiene av svindelpunktene bare havner i trening}
\alt{B}{Oversamplingen fjerner normale transaksjoner fra trening, så modellen undertilpasser og recall blir upålitelig}
\alt{C}{Kopier av de samme svindelpunktene havner både i trening og test, så modellen testes på punkter den har sett}
\alt{D}{Recall avhenger ikke av klassefordelingen, så resultatet er gyldig, mens accuracy ville blitt for høy}
\end{alternativer}
\makspoeng

\delstart{3}
\sporsmal{1}
I et binært klassifiseringsproblem har 20\,\% av datapunktene label $y=1$. En enkel modell predikerer
sannsynligheten $0.4$ for $y=1$ for alle punktene. La $p = (0.2,\ 0.8)$ være fordelingen til labelene og
$q = (0.4,\ 0.6)$ fordelingen modellen predikerer. Entropien er $H(p) = -\sum_i p_i \log_2 p_i$ og
cross-entropy er $H(p,q) = -\sum_i p_i \log_2 q_i$.

Hva blir $H(p)$ og $H(p,q)$?
\velg{Velg ett alternativ:}
\begin{alternativer}
\alt{A}{$H(p) = 0.72$ og $H(p,q) = 1.12$}
\alt{B}{$H(p) = 0.97$ og $H(p,q) = 0.85$}
\alt{C}{$H(p) = 0.50$ og $H(p,q) = 0.59$}
\alt{D}{$H(p) = 0.72$ og $H(p,q) = 0.85$}
\end{alternativer}
\sporsmal{2}
Notatene skiller mellom \textbf{parametrisk} og \textbf{ikke-parametrisk} modellering.

Hvilket utsagn er riktig?
\velg{Velg ett alternativ:}
\begin{alternativer}
\alt{A}{Beslutningstrær er ikke-parametriske fordi de ikke har hyperparametre som må velges før treningen}
\alt{B}{Logistisk regresjon er parametrisk fordi funksjonsformen er gitt, mens et tres struktur bestemmes under trening}
\alt{C}{Beslutningstrær er parametriske fordi terskelverdiene i splittene er parametre som læres fra dataene}
\alt{D}{Logistisk regresjon er ikke-parametrisk fordi antallet vekter øker med antallet datapunkter i treningssettet}
\end{alternativer}
\sporsmal{3}
Du vil estimere simultanfordelingen $p(\mathbf{x} \mid C_k)$ for én klasse uten Naive Bayes-antakelsen, ved hjelp
av et histogram. Du deler hver feature inn i 10 bins og vil ha minst 10 datapunkter i hver bin.

Hvor mange datapunkter trenger du (for denne klassen) hvis du har 6 features?
\velg{Velg ett alternativ:}
\begin{alternativer}
\alt{A}{$10^{7}$}
\alt{B}{$10^{6}$}
\alt{C}{$6^{10}$}
\alt{D}{$600$}
\end{alternativer}
\sporsmal{4}
Tabellen viser Titanic-treningsdataene fordelt på \emph{Sex} og \emph{Survived}.
\begin{center}
\begin{tabular}{lcc}
\toprule
 & \emph{Survived} $=0$ & \emph{Survived} $=1$ \\
\midrule
female & 81 & 233 \\
male & 468 & 109 \\
\bottomrule
\end{tabular}
\end{center}
En trestump splitter på \emph{Sex}, og hver løvnode predikerer klassen som dominerer i noden.

Hva blir treningsaccuracy for stumpen, og for en modell som alltid predikerer klassen som dominerer i hele datasettet?
\velg{Velg ett alternativ:}
\begin{alternativer}
\alt{A}{$0.79$ for stumpen og $0.38$ for majoritetsmodellen}
\alt{B}{$0.79$ for stumpen og $0.62$ for majoritetsmodellen}
\alt{C}{$0.67$ for stumpen og $0.62$ for majoritetsmodellen}
\alt{D}{$0.67$ for stumpen og $0.38$ for majoritetsmodellen}
\end{alternativer}
\makspoeng

\delstart{4}
\sporsmal{1}
Du trener $f(\mathbf{x}) = \beta_0 + \beta_1 x_1 + \beta_2 x_2$ med full-batch gradient descent på MSE.
Etter mange epoker er gradientvektoren
$\bigl(\partial L/\partial\beta_0,\ \partial L/\partial\beta_1,\ \partial L/\partial\beta_2\bigr)$ nesten lik nullvektoren.

Hva kan du konkludere?
\velg{Velg ett alternativ:}
\begin{alternativer}
\alt{A}{Treningstapet er nær null, så modellen gjengir treningsdataene nesten perfekt, men kan overtilpasse}
\alt{B}{Modellen har verken over- eller undertilpasset, men treningen bør fortsette til gradienten er null}
\alt{C}{Læringsraten er for lav, så parametrene står nesten stille, men en større læringsrate gir lavere tap}
\alt{D}{Parametrene ligger nær det globale minimumet av treningstapet, men tapet der kan fortsatt være høyt}
\end{alternativer}
\sporsmal{2}
Et regresjonstre vurderer splitten $x \le 4.5$ for dataene under. Variansen i en node beregnes som
gjennomsnittlig kvadrert avvik fra nodens gjennomsnitt (del på antall punkter i noden).
\begin{center}
\begin{tabular}{lcccccccc}
\toprule
$x$ & 1 & 2 & 3 & 4 & 5 & 6 & 7 & 8\\
$y$ & 3 & 4 & 3 & 5 & 9 & 10 & 8 & 9\\
\bottomrule
\end{tabular}
\end{center}
Hva blir den vektede variansen etter splitten, $\frac{n_1}{n}\text{Var}_1 + \frac{n_2}{n}\text{Var}_2$?
\velg{Velg ett alternativ:}
\begin{alternativer}
\alt{A}{$1.19$}
\alt{B}{$0.59$}
\alt{C}{$7.48$}
\alt{D}{$0.69$}
\end{alternativer}
\sporsmal{3}
Du har bare 120 datapunkter. To ulike tilfeldige 80/20-splitter av de samme dataene gir test-MSE 3.1 og 4.6
for den samme modelltypen og de samme hyperparametrene.

Hva er den beste tolkningen, og hva er et passende tiltak?
\velg{Velg ett alternativ:}
\begin{alternativer}
\alt{A}{Forskjellen skyldes den irredusible støyen, og den blir typisk mindre hvis testsettet gjøres enda mindre}
\alt{B}{Estimatet avhenger av hvilke 24 punkter som havner i testsettet; k-fold kryssvalidering gir et mer stabilt estimat}
\alt{C}{Den laveste verdien, 3.1, er det beste anslaget, siden den splitten ga de mest representative testdataene}
\alt{D}{Modellen overtilpasser, siden test-MSE varierer; den bør gjøres enklere til begge splittene gir omtrent like stor feil}
\end{alternativer}
\sporsmal{4}
En regresjonsmodell trenes på mange ulike treningssett trukket fra samme fordeling. I et bestemt punkt
$x$ er den sanne verdien $F(x) = 10$. Modellens prediksjoner i $x$ har gjennomsnitt 12 og varians 3 over
treningssettene, og den irredusible støyen har varians $\sigma^2 = 1$.

Hva er forventet kvadrert feil $\mathbb{E}_D\big[(y - f(x;D))^2\big]$ i punktet $x$?
\velg{Velg ett alternativ:}
\begin{alternativer}
\alt{A}{$8$}
\alt{B}{$6$}
\alt{C}{$7$}
\alt{D}{$5$}
\end{alternativer}
\makspoeng

\delstart{5}
\sporsmal{1}
Ensemble-metoder skiller mellom \textbf{trening} og \textbf{inferens}.

Hvilket utsagn er riktig?
\velg{Velg ett alternativ:}
\begin{alternativer}
\alt{A}{I boosting trenes modellene etter hverandre, og ved inferens kombineres prediksjonene fra alle modellene}
\alt{B}{I boosting trenes modellene etter hverandre, men ved inferens brukes typisk bare den siste modellen}
\alt{C}{I bagging trenes modellene etter hverandre, og ved inferens teller de senere modellene mest i snittet}
\alt{D}{I boosting kan modellene trenes parallelt, men ved inferens må de evalueres etter hverandre i fast rekkefølge}
\end{alternativer}
\sporsmal{2}
Du trener \textbf{gradient boosting} med trær av dybde 10 og 2000 boosting-runder. Treningsfeilen går mot 0,
mens valideringsfeilen er lavest etter rundt 300 runder og deretter øker jevnt.

Hva er den mest sannsynlige forklaringen, og hva er et fornuftig tiltak?
\velg{Velg ett alternativ:}
\begin{alternativer}
\alt{A}{Valideringsdataene har trolig en annen fordeling; feilen bør heller måles på testsettet underveis}
\alt{B}{Ensemblet undertilpasser fordi hvert tre er en svak lærer; flere runder og dypere trær vil hjelpe}
\alt{C}{Boosting reduserer først bias og deretter varians, så valideringsfeilen vil trolig synke igjen med flere runder}
\alt{D}{Ensemblet overtilpasser; antall runder og tredybde bør velges ut fra valideringsfeilen}
\end{alternativer}
\sporsmal{3}
I \textbf{AdaBoost} har fem datapunkter vekt 0.2 hver. En stump klassifiserer ett av punktene feil, så den
vektede feilraten er $\varepsilon = 0.2$ og $\alpha = \tfrac12 \ln\frac{1-\varepsilon}{\varepsilon}$. Vektene
oppdateres med $w_i \leftarrow w_i e^{\alpha}$ for feilklassifiserte punkter og $w_i \leftarrow w_i e^{-\alpha}$ for
riktig klassifiserte punkter, og normaliseres deretter slik at de summerer til 1.

Hva blir de nye vektene?
\velg{Velg ett alternativ:}
\begin{alternativer}
\alt{A}{Det feilklassifiserte punktet får 0.400, og hvert av de andre får 0.100}
\alt{B}{Det feilklassifiserte punktet får 0.333, og hvert av de andre får 0.167}
\alt{C}{Det feilklassifiserte punktet får 0.500, og hvert av de andre får 0.125}
\alt{D}{Det feilklassifiserte punktet får 0.059, og hvert av de andre får 0.235}
\end{alternativer}
\sporsmal{4}
I \textbf{AdaBoost} får hver svake lærer stemmevekten $\alpha = \tfrac12\ln\frac{1-\varepsilon}{\varepsilon}$,
der $\varepsilon$ er den vektede feilraten. Tre stumper har $\varepsilon = 0.10$, $\varepsilon = 0.30$ og $\varepsilon = 0.45$.

Hvilket utsagn om stemmevektene er riktig?
\velg{Velg ett alternativ:}
\begin{alternativer}
\alt{A}{Stumpen med $\varepsilon = 0.10$ får størst stemmevekt, og en stump med $\varepsilon = 0.60$ ville fått vekt 0}
\alt{B}{Alle tre får lik stemmevekt, siden AdaBoost kombinerer de svake lærerne ved vanlig flertallsavstemning}
\alt{C}{Stumpen med $\varepsilon = 0.45$ får størst stemmevekt, siden den er trent på de vanskeligste punktene}
\alt{D}{Stumpen med $\varepsilon = 0.10$ får størst stemmevekt, og den med $\varepsilon = 0.45$ får en vekt nær 0}
\end{alternativer}
\makspoeng

\delstart{6}
\sporsmal{1}
Et nettverk har et skjult lag og tre output-noder. For et datapunkt vil én output-node få lavere tap om
aktiveringen $a_k$ til node $k$ i det skjulte laget øker, mens en annen output-node vil få lavere tap om
$a_k$ minker.

Hvordan håndterer backpropagation dette når gradienten for vektene inn til node $k$ skal beregnes?
\velg{Velg ett alternativ:}
\begin{alternativer}
\alt{A}{Bidragene fra alle output-nodene legges sammen, og den samlede ønskede endringen sendes videre bakover}
\alt{B}{Output-noden med størst feil bestemmer endringen alene, og bidragene fra de andre nodene ignoreres her}
\alt{C}{Aktiveringen $a_k$ settes direkte til verdien som flertallet av output-nodene ønsker, før vektene oppdateres}
\alt{D}{Vektene inn til node $k$ oppdateres én gang per output-node, i fast rekkefølge fra første til siste node}
\end{alternativer}
\sporsmal{2}
Se på følgende kode:
\begin{lstlisting}[language=Python]
model = keras.Sequential()
model.add(keras.Input(shape=(3,)))
model.add(Dense(6, activation="relu"))
x = ops.ones((10, 4))
out = model(x)
\end{lstlisting}
Hva skjer når den siste linjen kjøres?
\velg{Velg ett alternativ:}
\begin{alternativer}
\alt{A}{Laget bygges på nytt for 4 features, og \texttt{out} får shape \texttt{(10, 6)} med nye vekter}
\alt{B}{Kallet gir en feilmelding, siden laget er bygget for nøyaktig 3 input-features}
\alt{C}{Kallet lykkes, og \texttt{out} får shape \texttt{(4, 6)}, siden hver feature gir én rad i output}
\alt{D}{Kallet lykkes, og \texttt{out} får shape \texttt{(10, 6)}, siden den fjerde featuren ignoreres}
\end{alternativer}
\sporsmal{3}
Et lite nettverk bruker Heaviside-funksjonen $H(z) = 1$ hvis $z \ge 0$, ellers 0:
\[
h_1 = H(x_1 + x_2 - 0.5), \qquad h_2 = H(-x_1 - x_2 + 1.5), \qquad y = H(h_1 + h_2 - 1.5).
\]
For hvilke binære input $(x_1, x_2)$ blir $y = 1$?
\velg{Velg ett alternativ:}
\begin{alternativer}
\alt{A}{Bare for $(0,1)$ og $(1,0)$}
\alt{B}{For alle unntatt $(0,0)$}
\alt{C}{Bare for $(1,1)$}
\alt{D}{Bare for $(0,0)$ og $(1,1)$}
\end{alternativer}
\sporsmal{4}
Se på følgende kode:
\begin{lstlisting}[language=Python]
model = keras.Sequential([keras.Input(shape=(20,)),
                          Dense(32, activation="relu"),
                          Dense(5, activation="softmax")])
W, b = model.layers[0].get_weights()
print(W.shape, b.shape)
\end{lstlisting}
Hva skrives ut?
\velg{Velg ett alternativ:}
\begin{alternativer}
\alt{A}{\texttt{(32, 5) (5,)}}
\alt{B}{\texttt{(20, 32) (32,)}}
\alt{C}{\texttt{(20, 32) (20,)}}
\alt{D}{\texttt{(32, 20) (32,)}}
\end{alternativer}
\makspoeng

\delstart{7}
\sporsmal{1}
Du sammenligner \textbf{DBSCAN} med k-means på et datasett med støy.

Hvilket utsagn om DBSCAN er riktig?
\velg{Velg ett alternativ:}
\begin{alternativer}
\alt{A}{Antall klynger må oppgis, og punkter som ligger langt fra alle sentroidene får label $-1$}
\alt{B}{Antall klynger må ikke oppgis, men alle punktene blir tilordnet den nærmeste klyngen}
\alt{C}{Antall klynger må ikke oppgis, og punkter som ikke kan nås fra noe kjernepunkt får label $-1$}
\alt{D}{Antall klynger må oppgis, men støypunktene samles i en egen klynge med label 0}
\end{alternativer}
\sporsmal{2}
En clustering-algoritme beskriver hver klynge med en middelvektor og en kovariansmatrise, som estimeres
iterativt fra dataene. Klyngene kan dermed ha ulik størrelse og elliptisk form, og hvert punkt får en grad av
tilhørighet til hver av klyngene i stedet for én fast klynge.

Hvilken av notatenes hovedkategorier hører algoritmen til?
\velg{Velg ett alternativ:}
\begin{alternativer}
\alt{A}{Sentroidebasert, siden middelvektoren fungerer som en sentroide som punktene tilordnes iterativt}
\alt{B}{Fordelingsbasert, siden middelvektoren og kovariansmatrisen er parametrene til en Gauss per klynge}
\alt{C}{Hierarkibasert, siden hvert punkt tilhører flere klynger samtidig, med ulik grad av tilhørighet}
\alt{D}{Tetthetsbasert, siden kovariansmatrisen beskriver hvor tett punktene ligger rundt middelvektoren}
\end{alternativer}
\sporsmal{3}
Et stort datasett består av 4 like store, godt separerte klynger med svært mange punkter i hver. Du kjører
k-means med $k = 4$ og plasserer de fire startsentroidene på tilfeldig valgte datapunkter.

Hva er (tilnærmet) sannsynligheten for at hver klynge får nøyaktig én startsentroide?
\velg{Velg ett alternativ:}
\begin{alternativer}
\alt{A}{$0.004$}
\alt{B}{$0.042$}
\alt{C}{$0.016$}
\alt{D}{$0.094$}
\end{alternativer}
\sporsmal{4}
En butikk klynger 2 millioner kunder, hver beskrevet med 12 standardiserte features, med k-means
($k = 6$). Modellen skal senere brukes til å plassere nye kunder i en av klyngene.

Hvor mange tall må den trente modellen lagre for dette, når vi ser bort fra $k$ og skaleringsparametrene?
\velg{Velg ett alternativ:}
\begin{alternativer}
\alt{A}{$6$}
\alt{B}{$24\,000\,000$}
\alt{C}{$72$}
\alt{D}{$18$}
\end{alternativer}
\makspoeng

\delstart{8}
\sporsmal{1}
Et overvåkingssystem kjører følgende hver dag på dagens transaksjoner \texttt{X\_dag}:
\begin{lstlisting}[language=Python]
distances = kmeans.transform(X_dag).min(axis=1)
threshold = np.percentile(distances, 95)
flagget = distances > threshold
\end{lstlisting}
En dag består alle transaksjonene av helt vanlige kjøp, uten noen svindel. Hva skjer?
\velg{Velg ett alternativ:}
\begin{alternativer}
\alt{A}{Ingen transaksjoner flagges, siden ingen avstander er uvanlig store når alle kjøpene er normale}
\alt{B}{Rundt 95\,\% av transaksjonene flagges, siden 95-persentilen gjør at de fleste havner over terskelen}
\alt{C}{Antall flaggede avhenger av $k$, siden hver klynge typisk får sin egen terskel på 5\,\% av punktene}
\alt{D}{Rundt 5\,\% av transaksjonene flagges likevel, siden terskelen er beregnet ut fra dagens egne avstander}
\end{alternativer}
\sporsmal{2}
Et sentrert datasett med to features har første hovedkomponent $\text{PC1} = (0.6,\ 0.8)$. Du reduserer
punktet $\mathbf{x} = (3,\ 1)$ til én dimensjon med PCA, og transformerer det deretter tilbake til det
opprinnelige rommet (\texttt{inverse\_transform}).

Hvilket punkt får du tilbake?
\velg{Velg ett alternativ:}
\begin{alternativer}
\alt{A}{$(2.60,\ 0.00)$}
\alt{B}{$(3.00,\ 1.00)$}
\alt{C}{$(1.56,\ 2.08)$}
\alt{D}{$(1.80,\ 0.80)$}
\end{alternativer}
\sporsmal{3}
Du bruker PCA til å redusere S-formede data i tre dimensjoner (fra \texttt{make\_s\_curve}) til to
dimensjoner, og resultatet bevarer strukturen godt.

Hva er den beste forklaringen?
\velg{Velg ett alternativ:}
\begin{alternativer}
\alt{A}{PCA er ikke-lineær og kan rette ut S-formen, slik t-SNE gjør med ikke-lineære strukturer}
\alt{B}{Dataene varierer lite i én retning, så to lineære akser fanger nesten all variansen}
\alt{C}{PCA bruker fargene (labels) til å orientere komponentene, så klassene holdes adskilt}
\alt{D}{PCA fjerner retningen med størst varians, siden den typisk består mest av støy}
\end{alternativer}
\sporsmal{4}
Se på følgende kode:
\begin{lstlisting}[language=Python]
pca = PCA().fit(X_train)
r = pca.explained_variance_ratio_     # [0.55, 0.25, 0.12, 0.05, 0.03]
k = np.argmax(np.cumsum(r) >= 0.90) + 1
print(k)
\end{lstlisting}
Hva skrives ut?
\velg{Velg ett alternativ:}
\begin{alternativer}
\alt{A}{\texttt{3}}
\alt{B}{\texttt{4}}
\alt{C}{\texttt{5}}
\alt{D}{\texttt{2}}
\end{alternativer}
\makspoeng

\delstart{9}
\sporsmal{1}
Et \textbf{DQN} trenes med både et Q-nettverk og et \textbf{target network}.

Hva er rollen til target network?
\velg{Velg ett alternativ:}
\begin{alternativer}
\alt{A}{Det beregner rewarden $R_{t+1}$, siden miljøet bare gir reward når en episode er over}
\alt{B}{Det velger handlingene under treningen, mens Q-nettverket bare brukes etter at treningen er ferdig}
\alt{C}{Det gir stabile mål: $\max_A Q(S_{t+1}, A)$ beregnes med en kopi som bare oppdateres periodisk}
\alt{D}{Det lagrer tidligere overganger, slik at de kan trekkes tilfeldig i batcher under treningen}
\end{alternativer}
\sporsmal{2}
I DQN lagres overganger $(S_t, A_t, R_{t+1}, S_{t+1}, \text{done})$ i en \textbf{replay buffer}, og
Q-nettverket trenes på tilfeldig trukne batcher fra bufferen.

Hva er hovedgevinsten sammenlignet med å trene på hver overgang idet den skjer?
\velg{Velg ett alternativ:}
\begin{alternativer}
\alt{A}{Påfølgende steg havner sjelden i samme batch, og hver erfaring kan brukes i flere oppdateringer}
\alt{B}{Tilfeldig trekking fra bufferen gir nok utforskning, så $\varepsilon$ kan settes til 0 tidligere i treningen}
\alt{C}{Overgangene fra én episode trenes samlet, slik at rewarden ved målet raskere forplanter seg bakover}
\alt{D}{Målene $R_{t+1} + \gamma\max_A Q(S_{t+1},A)$ holdes faste mellom oppdateringene, så treningen blir stabil}
\end{alternativer}
\sporsmal{3}
En agent bruker Q-læring med $Q_{new}(S_t,A_t) \leftarrow (1-\alpha)\,Q(S_t,A_t) + \alpha\,\big(R_{t+1} + \gamma \max_A Q(S_{t+1},A)\big)$,
med $\alpha = 0.4$ og $\gamma = 0.9$. Agenten har $Q(S_t, A_t) = 2.0$, mottar $R_{t+1} = 1$, og Q-verdiene i den nye
tilstanden er $Q(S_{t+1}, \cdot) = (3.0,\ 4.0,\ 1.5)$.

Hva blir $Q_{new}(S_t, A_t)$?
\velg{Velg ett alternativ:}
\begin{alternativer}
\alt{A}{$3.20$}
\alt{B}{$4.60$}
\alt{C}{$3.56$}
\alt{D}{$3.04$}
\end{alternativer}
\sporsmal{4}
Se på Q-læringsregelen $Q_{new}(S_t,A_t) \leftarrow (1-\alpha)\,Q(S_t,A_t) + \alpha\,\big(R_{t+1} + \gamma \max_A Q(S_{t+1},A)\big)$.

Hva skjer med $\alpha = 0$ og med $\alpha = 1$?
\velg{Velg ett alternativ:}
\begin{alternativer}
\alt{A}{Med $\alpha = 0$ teller bare umiddelbar reward; med $\alpha = 1$ vektes fremtidige rewards like mye som nåværende}
\alt{B}{Med $\alpha = 0$ endres Q-verdiene aldri; med $\alpha = 1$ erstattes den gamle verdien helt av det nye estimatet}
\alt{C}{Med $\alpha = 0$ konvergerer Q-verdiene raskest; med $\alpha = 1$ blir stegene så store at Q-verdiene divergerer}
\alt{D}{Med $\alpha = 0$ velges alltid den grådige handlingen; med $\alpha = 1$ velges alle handlingene helt tilfeldig}
\end{alternativer}
\makspoeng

\delstart{10}
\sporsmal{1}
Tabellen viser SHAP-verdiene for tre features på fire datapunkter.
\begin{center}
\begin{tabular}{lcccc}
\toprule
Feature & Punkt 1 & Punkt 2 & Punkt 3 & Punkt 4\\
\midrule
A & $+0.30$ & $-0.30$ & $+0.25$ & $-0.25$\\
B & $+0.10$ & $+0.12$ & $+0.08$ & $+0.10$\\
C & $-0.20$ & $+0.05$ & $-0.15$ & $+0.35$\\
\bottomrule
\end{tabular}
\end{center}
Hvilken feature havner øverst i et beeswarm-plott av disse SHAP-verdiene?
\velg{Velg ett alternativ:}
\begin{alternativer}
\alt{A}{B, siden den har størst gjennomsnittlig SHAP-verdi}
\alt{B}{C, siden den har den største enkeltstående SHAP-verdien av alle}
\alt{C}{A, siden den har størst gjennomsnittlig absolutt SHAP-verdi}
\alt{D}{Ingen, siden lokale SHAP-verdier ikke kan slås sammen}
\end{alternativer}
\sporsmal{2}
En lånemodell forklares med SHAP. Bakgrunnsdatasettet byttes fra et tilfeldig utvalg av alle søkere til et
utvalg av bare søkere som fikk innvilget lån. Prediksjonen $f(\mathbf{x}^*)$ for søkeren som forklares, er den samme.

Hvilket utsagn er riktig?
\velg{Velg ett alternativ:}
\begin{alternativer}
\alt{A}{SHAP-verdiene endres typisk, men $\phi_0$ er uendret, siden $\phi_0$ er det samme som modellens bias $b$}
\alt{B}{SHAP-verdiene er uendret mens $\phi_0$ endres, så $\phi_0 + \sum_j \phi_j$ blir ikke lenger lik $f(\mathbf{x}^*)$}
\alt{C}{Både $\phi_0$ og SHAP-verdiene kan endres, men $\phi_0 + \sum_j \phi_j$ er fortsatt lik $f(\mathbf{x}^*)$}
\alt{D}{$\phi_0$ og SHAP-verdiene er uendret, siden de bare avhenger av modellen og datapunktet $\mathbf{x}^*$}
\end{alternativer}
\sporsmal{3}
Se på følgende kode:
\begin{lstlisting}[language=Python]
df = pd.DataFrame({"Pclass": [1, 1, 3, 3, 3, 3, 2, 2],
                   "Sex":    ["f", "m", "m", "m", "f", "m", "f", "m"]})
p = pd.crosstab(df.Pclass, df.Sex, normalize="index")
print(p.loc[3, "m"])
\end{lstlisting}
Hva skrives ut?
\velg{Velg ett alternativ:}
\begin{alternativer}
\alt{A}{\texttt{0.75}}
\alt{B}{\texttt{0.6}}
\alt{C}{\texttt{3}}
\alt{D}{\texttt{0.375}}
\end{alternativer}
\sporsmal{4}
Se på følgende kode:
\begin{lstlisting}[language=Python]
y_true = [1, 0, 1, 1]
y_pred = [1, 0, 0, 1]
print(np.mean(y_true == y_pred), np.mean(np.array(y_true) == np.array(y_pred)))
\end{lstlisting}
Hva skrives ut?
\velg{Velg ett alternativ:}
\begin{alternativer}
\alt{A}{\texttt{1.0 0.75}}
\alt{B}{\texttt{0.0 0.75}}
\alt{C}{\texttt{0.75 0.75}}
\alt{D}{\texttt{0.75 0.0}}
\end{alternativer}
\makspoeng

\clearpage
{\sffamily\bfseries\fontsize{18}{22}\selectfont Løsningsforslag}\par\vspace{0.3em}
\noindent\rule{\textwidth}{0.4pt}\par\vspace{1.2em}
Riktig alternativ per spørsmål (bokstaven står i rundingen foran hvert alternativ).\par\vspace{1em}
\begin{center}\sffamily\renewcommand{\arraystretch}{1.35}
\begin{tabular}{l*{4}{>{\centering\arraybackslash}p{1.6cm}}}
\toprule
& \textbf{Sp.\ 1} & \textbf{Sp.\ 2} & \textbf{Sp.\ 3} & \textbf{Sp.\ 4}\\
\midrule
Part 1 & D & A & C & D\\
Part 2 & C & D & A & C\\
Part 3 & D & B & A & B\\
Part 4 & D & B & B & A\\
Part 5 & A & D & C & D\\
Part 6 & A & B & A & B\\
Part 7 & C & B & D & C\\
Part 8 & D & C & B & A\\
Part 9 & C & A & D & B\\
Part 10 & C & C & A & B\\
\bottomrule
\end{tabular}
\end{center}
\end{document}
